\section{Introduction} \label{sec:intro}

Partial Label Learning (PLL) addresses the challenge of weakly supervised learning where each training instance is annotated with a set of candidate labels, only one of which is the ground truth~\cite{cour2009learning, zhang2015solving, feng2020provably}. While standard PLL assumes the ground truth is strictly contained within the candidate set, real-world scenarios such as web-scale crawling and crowdsourced annotation inevitably introduce noise. This gives rise to the more general and challenging problem of Noisy Partial Label Learning (NPLL)~\cite{wang2022pico+}, where the true label may be absent from the candidate set, rendering conventional disambiguation assumptions invalid.

The candidate set is not necessarily random noise. Instance-dependent PLL shows that incorrect candidate labels are often related to the instance rather than sampled uniformly. VALEN models such feature-dependent candidate labels through latent label distributions~\cite{xu2021idpll}, while IDGP decomposes the annotation process into the generation of the true label followed by feature-related incorrect candidates~\cite{qiao2023idgp}. POP further updates the model and progressively purifies each candidate set in instance-dependent PLL~\cite{xu2022pop}. These studies indicate that the observed candidate set can contain structured weak-supervision information induced by the instance and the annotation process. In NPLL and crowdsourced annotation, however, this information is also noisy and may miss the ground-truth class. A robust method should therefore exploit the observed prior without assuming that it is always correct.

To mitigate the impact of label noise, recent approaches have integrated contrastive learning~\cite{wang2022pico, wang2022pico+}, sample separation~\cite{lian2023irnet, Shi2023Unreliable}, candidate-set refinement~\cite{qiao2023fredis}, candidate-label reconstruction~\cite{peng2025noise}, or label-importance adjustment~\cite{xu2023alim} into the PLL framework. A representative graph-based method, PALS~\cite{pals2025}, leverages weighted nearest-neighbor pseudo-labelling and label smoothing to construct reliable image--label pairs. These methods demonstrate that geometry, model prediction, and the original candidate set are all useful sources of evidence. They also reveal a design tension: training supervision is often made increasingly model-dependent through purification, refinement, correction, reconstruction, pseudo-target updating, sample promotion, or partial-label augmentation. We call any supervision object inherited by later epochs a \emph{persistent supervision state}; this includes candidate membership, reconstructed candidate sets, soft pseudo-targets, label-confidence or importance weights, and reliable/unreliable promotion states. Such operations are motivated by the belief that reducing candidate noise should make the supervision cleaner. However, when model-induced decisions are carried across epochs as persistent supervision state, early prediction errors can remain in the working supervision used later in training.

We challenge the necessity of this persistent carry-over route. The observed candidate set is noisy, but it is not supervision-free. It contains structured information induced by instance ambiguity, annotator uncertainty, and crowd disagreement. Instead of asking only how to purify the candidate set or how to inherit a corrected supervision state, we ask how to extract reliable supervision from the native weak prior without writing epoch-local model decisions back into any later-epoch supervision state. This work argues that NPLL should not be viewed only as a problem of candidate-set correction: high-performance NPLL can be obtained by restoring the native weak-supervision record at each epoch and extracting low-error active supervision from the restored prior.

In this work, we propose \textbf{NSE} (\textbf{N}on-Destructive \textbf{S}upervision \textbf{E}xtraction), a framework designed around a reset-before-extraction principle. NSE does not claim that weak supervision should never be transformed. Instead, it separates the native weak-supervision record from epoch-local working supervision and from model-side learning state. At the beginning of each epoch, the working prior is restored from the original candidate/crowd prior; Topology-DAES propagation, model-evidence projection, graph-semantic fusion, reliable selection, salvage, and pseudo-label training then operate on temporary evidence that is discarded after the epoch. The term non-destructive is therefore restricted to persistent supervision state: model parameters, optimizer state, prototypes, and diagnostic histories are still learned across epochs, but model-induced candidate membership, pseudo-targets, or sample-promotion decisions are not inherited as the supervision source for later epochs. NSE addresses supervision extraction with \textbf{Topology-DAES}, a topology-guided Dynamic Adaptive Entropy Similarity kernel. Consistent neighborhoods provide sharp geometric evidence, while ambiguous weak-label neighborhoods are prevented from aggressively spreading unreliable labels. Model predictions are used only as temporary semantic evidence and must pass through the native prior before entering the second Topology-DAES propagation. Training is then performed on an active set consisting of reliable pseudo-labels and samples recovered by stable model--KNN--prototype consensus, while ordinary non-reliable samples are assigned zero sampling weight. Our contributions are summarized as follows:

\begin{itemize}
    \item \textbf{We identify a persistent supervision-state risk in NPLL.} Repeatedly purifying, correcting, reconstructing, promoting, or reweighting weak labels can make later-epoch supervision model-dependent and may turn early biased predictions into inherited supervision state.

    \item \textbf{We propose epoch-wise source restoration.} NSE restores the working prior from the native candidate/crowd record at the beginning of each epoch and focuses on extracting reliable active supervision through graph-semantic evidence fusion.

    \item \textbf{We introduce Topology-DAES as the supervision extraction kernel of NSE.} It uses masked-entropy topology reliability to measure local agreement between neighbor votes and the native weak prior, and uses entropy-adaptive temperature to control propagation under ambiguity.

    \item \textbf{We combine Topology-DAES geometric evidence with prior-constrained model evidence.} Model predictions provide temporary semantic evidence and must be projected by the native prior before the second propagation stage.

    \item \textbf{We introduce active-set training with consensus-based salvage.} Reliable selection and salvage update training eligibility, not candidate membership.

    \item \textbf{We empirically challenge the necessity of persistent supervision-state carry-over.} Across synthetic, hierarchical, and crowdsourced NPLL settings, and through prior-aligned persistent-state proxy tests, high performance can be achieved without inheriting model-induced supervision-state modifications across epochs.
\end{itemize}

